# Profil J1 — Prise de notes brute (Stagiaire, 5 mois)

**Code :** J1  
**Fonction :** Stagiaire  
**Ancienneté :** 5 mois  
**Contexte :** Wavestone Luxembourg. Propales, institutions européennes, chef de projet, support livrables.

---

## T0 — Cadrage

**Q : Parcours, comment arrivé ici ?**
- 5 mois... cherchais exp en conseil. Wavestone Lux recrutait. Un peu par hasard.
- Fais bcp de propales. Support sur missions instits européennes. Chef de projet aussi, interface seniors / équipes tech.
- Part d'écrit ? Énorme. 70%. PPT, notes cadrage, CR. Le reste 30% oral : points manager, calls clients.

---

## T1 — Usage réel

**Q : Dernière fois utilisé outil génératif ?**
- Semaine dernière. Propale. Dossier à rendre. 10 slides sur sujet que maîtrisais moyennement.
- Copilot direct dans PPT. Pas ChatGPT. Demandé structure titres/sous-parties sur 3 docs.
- 5 min... contre 1h à tout relire. Après, j'ai repris chaque slide. Contenu générique mais squelette sympa. J'ai mis la viande autour.

**Q : Sur quelles tâches au quotidien ?**
- Tout le temps : reformulation. Notes réunion fouillis -> "rends-moi structuré". Là-dessus c'est un tueur.
- Traductions aussi. Anglais/français.
- Jamais sur slides avec données chiffrées confidentielles. Dealbreaker.

**Q : Comment appris ?**
- Seul. Pas de formation. Webinaire récent mais j'étais en mission. Testé, trifouillé.
- Suivi progrès : mars quasi rien. Texte brut moche. Mai actions sans styling (transitions). Récemment diapos entières en un prompt. Changé en 3 mois.

**Q : Temps gagné mis sur quoi ?**
- Propale, gagné 2h. Utilisées pour relire 3x le doc. Vérifier ancrage client.
- Permet de passer + de temps sur qualité fond, moins forme/squelette. Contre-intuitif, mais fait mieux sur l'essentiel.

**Q : Situation où plus d'effort que de le faire soi-même ?**
- Oui. Sujet hyper technique, règlementation bancaire. Texte vague, presque juste, refs à directives inexistantes.
- 20min à reformuler prompts... rien. Réécrit 15min. Depuis, trop pointu sans sources, je fais direct.

**Q : Vous en parlez entre vous ?**
- Entre stagiaires/juniors, oui. Prompts, astuces.
- Avec seniors, plus discret. Pas dire "j'ai utilisé Copilot" au manager. Non-dit. Tout le monde l'utilise, personne le dit. Peur qu'on nous dise "tu réfléchis plus".

---

## T2 — Apprentissage

**Q : Comment on apprend le métier aujourd'hui ?**
- Différent. Notes/CR plus vite, mais pas sûr d'apprendre la même chose.
- Avant : trier, reformuler. Maintenant : colle notes, "structure-moi". Gagne temps, mais n'exerce pas synthèse. Plus "relecteur" que "rédacteur".

**Q : Par quelles tâches on apprend vraiment ?**
- Relecture manager. 1ère propale, 2j à faire. Manager a renvoyé plein de commentaires, titres, structure. Repris en comprenant *pourquoi*.
- Maintenant Copilot => version propre vite. Manager passe moins de temps sur forme, plus sur fond. Comms + intéressants ("raisonnement léger" vs "titres pas alignés"). Paradoxalement améliore relecture senior.
- Mais crainte : perdre capacité à faire version seule.

**Q : Tâches ingrates — CV par ex ?**
- Adaptation CV. Nuance, pertinence.
- Au début à la main. Puis Copilot : CV original + desc poste. Marchait bien. Une fois, dérapage. IA inventé projets inexistants (ou d'autres personnes). Si pas vérifié ligne à ligne, mensonge pur.
- Depuis, plus confiance sur les faits.

---

## T3 — Client et valeur

**Q : Confiance client repose sur quoi ?**
- Parler leur langage (métier, contraintes). + Responsabilité : "ces gens-là ont vérifié, engagés".
- Client ne demande pas IA, demande vrai. Confiance = nom sur le document.

**Q : Diriez-vous au client qu'une partie est IA ?**
- Non. Pas mentir si demande, mais pas mettre en avant.
- Risque : "pourquoi je paie un cabinet ?". Clients soupçonnent. On élude "technologies qui aident". Jamais dit "on génère des slides avec Copilot". Sujet tabou.

**Q : Livrable moins crédible ?**
- Erreurs factuelles. Chiffre faux => tout remis en cause. IA pas de statut. Erreur humaine excusable, erreur machine = faute pro.
- Donc très prudent sur chiffres/factuel. Vérifie tout. Si génère slide données, repasse avec sources.

**Q : Dans 10 ans, client paiera pour quoi ?**
- Capacité à prendre des risques. Dire vérité inconfortable. IA fait analyse/synthèse/mise en forme, pas prendre position. Pas dire "vous vous trompez". Payé pour jugement.

---

## T4 — Gouvernance

**Q : Règles sur usage ?**
- Officiellement, oui. Session y a 2 mois. Copilot approuvé, attention données confidentielles, anonymiser.
- Claude/ChatGPT = flou total. "Faites attention", pas de règle écrite.
- Utilise officieusement, très attention. Pas sûr que tout le monde prudent.

**Q : Avant livrable client ?**
- Revue hiérarchique. Version préliminaire -> manager -> commentaires -> corrige -> valid.
- Pas de "check IA". Personne demande "généré par ?". Justifier choix/chiffres/sources. Contrôle fond, pas processus. Managers soupçonnent mais ne posent pas.

**Q : Incident CV, qu'est-ce qui s'est passé ?**
- Dit à collègue stagiaire, rigolé. Pas dit manager. Honte.
- Pensait qu'elle allait dire pas fait attention. Repris CV main, envoyé.
- Frayeur. Si parti sans vérifier, mensonge client. Responsabilité sur moi.

**Q : Règle idéale ?**
- "Utiliser IA pour rédactionnel/structure/synthèse. Mais justifier chaque fait. Si pas justifiable, pas utiliser."
- Pas interdiction stricte, mais obligation vérification (sources, références, noms).
- Ce qui m'a sauvé = avoir gardé source de vérité à côté.

---

## T5 — Clôture

**Q : À la place de la direction, première décision ?**
- Clarifier statut outils non-officiels. Interdire ou approuver avec garde-fous.
- Flou dangereux. Formaliser "usage autorisé" / "interdit".
- Formation obligatoire pas sur prompts, mais sur détection hallucinations. Risque = consultants ne voient pas erreurs.

**Q : Aspect important non abordé ?**
- Temps psychologique. Gagne temps, mais pression.
- Slide en 30 sec => moins enclin à réfléchir fond avant. Lance outil, modifie. Plus vite, mais impression de ne pas maîtriser sujet. Illusion de compétence.
- Dans 6 mois, saurai faire structure seul ?

**Q : Point de vue utile à recueillir ?**
- Manager 10 ans. Voir comment il vit juniors avec livrables propres mais moins réfléchis.
- Lui valide, signe. Lui voit risque, moi accélérateur. Sa capacité à différencier livrable junior compétent vs junior IA.